%------------------------------------------------------------------------------%
\section{Funções Inversas}

Retomaremos ao Exemplo~\ref{ex:funcao_exp_1}, do capítulo anterior,
para introduzir a ideia de função Inversa e
posteriormente definiremos a função logarítmica que é a inversa da função
exponencial.

%------------------------------------------------------------------------------%
\begin{example}
  Considerando a situação descrita no Exemplo~\ref{ex:funcao_exp_1},
  se Heloísa não saldar parte de sua dívida, em que momento ela deverá o
  dobro do valor que pegou emprestado?
\end{example}
%------------------------------------------------------------------------------%

Já vimos que o valor a ser pago por Heloísa, $f(x)$, após $x$ meses que ela
contraiu o empréstimo é dado por
\[
  1000 \cdot 1,06^{x}.
\]
Essa função nos permite calcular o valor da imagem $f(x)$, dado que temos o
valor do domínio $x$. Observe que o problema proposto nos pergunta qual o valor
$x$ do domínio que é levado no valor $f(x)=2000$ da imagem.  Nosso problema
será facilmente resolvido se tivermos uma função que dado o valor de $f(x)$ ela
nos da o valor de $x$. Essa função é chamada função inversa, a qual passaremos
a definir agora.

%------------------------------------------------------------------------------%
\begin{figure}
  \begin{subfigure}[b]{0.5\linewidth}
    \centering
    \includegraphics[width=0.9\textwidth,trim={1cm 1cm 1cm 1cm},clip]{fun1}
    \caption{Função $f(x)$.}
    \label{fig:funcao_inv_f}
  \end{subfigure}%
  \begin{subfigure}[b]{0.5\linewidth}
    \centering
    \includegraphics[width=0.9\textwidth,trim={1cm 1cm 1cm 1cm},clip]{fun3}
    \caption{Função $h(x)$.}
    \label{fig:funcao_inv_h}
  \end{subfigure}
  \caption{Exemplos de funções.}
  \label{fig:funcoes_inv}
\end{figure}
%------------------------------------------------------------------------------%

A primeira pergunta a ser feita é: será que todas as funções possuem inversas?
Vamos observar duas situações em que as funções não possuem função inversa.

Na função $f\colon A\to B$, da Figura~\ref{fig:funcao_inv_f},
não existe $x\in A$ tal que $f(x)=a$.
Logo, não é possível definir uma função inversa já que não
teríamos como definir $f^{-1}(a)$.

Na função $h\colon C\to D$, da Figura~\ref{fig:funcao_inv_f}, $h(3)=b=h(4)$.
Logo, não é possível definir uma função inversa já que tem dois possíveis
valores para $h^{-1}(b)$.

Para definirmos a inversa de uma função $f$, nenhuma das situação acima podem
acontecer, uma função em que não ocorre essas situações é uma função bijetora,
cuja definição formal é dada por

%------------------------------------------------------------------------------%
\begin{definicao}{Função Bijetora}{def:bijetora}
  Uma função $f\colon A\to B$ é dita \emph{bijetora}\index{Função!bijetora}
  se satisfaz  as seguintes condições:
  \begin{enumerate}
    \item quaisquer elementos distintos de $A$ são levados por $f$ em elementos
      distintos de $B$, isto é, se para todo
      \[
       x, y\in A \text{ tal que } x \neq y, \text{ tem-se que } f(x) \neq f(y);
      \]

    \item se todo elemento $y \in B$ possui um representante
      \[
        x \in A \text{ tal que } f(x)=y.
      \]
  \end{enumerate}
\end{definicao}
%------------------------------------------------------------------------------%

Ou seja, $f$ é bijetora se para todo $y\in B$ existe um
\emph{único} $x\in A$ tal que $f(x)=y$.

Uma função $f\colon A\to B$ que satisfaz o Item $1$ da Definição
\ref{def:bijetora} é dita \emph{injetora}\index{Função!injetora}.

Uma função $f\colon A\to B$ que satisfaz o Item $2$ da Definição
\ref{def:bijetora} é dita \emph{sobrejetora}\index{Função!sobrejetora}.

Quando uma função $f\colon A\to B$ é bijetora, é possível definir uma
função $g\colon B \to A$ que, de certa forma, desfaz o que a $f$ faz.
Como para cada $y\in B$ existe um único $x\in A$ tal que $f(x)=y$,
definimos  $g(y)=x$ em que $f(x)=y$. Essa função assim definida é
chamada \emph{função inversa}\index{Função!inversa}
de $f$ e denotada por $f^{-1}$. Segue que
\[
  f^{-1}(y) = x
  \quad\Leftrightarrow\quad
  f(x) = y.
\]

Veja que pela definição de função, para que a função inversa esteja bem
definida, é necessário que $f\colon A\to B$ seja bijetiva. Pois nesse caso, a
sobrejetividade de $f$ garante que $f^{-1}(y)$ está definida para todo $y\in
B$ e a injetividade garante que o valor $f^{-1}(y)$ é único.

Note que se consideramos a função $f$ definida na sua imagem, ou seja
$f\colon D(A) \to Im(A)$, então para que $f$ possua inversa, é necessário
apenas que $f$ seja injetora.

%------------------------------------------------------------------------------%
\begin{figure}
  \begin{subfigure}[b]{0.5\linewidth}
    \centering
    \includegraphics[width=0.9\textwidth,trim={1cm 1cm 1cm 1cm},clip]{fun2}
    \caption{Função $g(x)$.}
    \label{fig:funcao_inv_g}
  \end{subfigure}%
  \begin{subfigure}[b]{0.5\linewidth}
    \centering
    \includegraphics[width=0.9\textwidth,trim={1cm 1cm 1cm 1cm},clip]{fun6}
    \caption{Função $g^{-1}(x)$.}
    \label{fig:funcao_inv_g_inv}
  \end{subfigure}
  \caption{Uma função e sua inversa.}
  \label{fig:funcao_g_e_inv_g}
\end{figure}
%------------------------------------------------------------------------------%

A função $g\colon C\to B$, da Figura~\ref{fig:funcao_inv_g},
é bijetora e sua inversa $g^{-1}\colon B\to C$,
Figura~\ref{fig:funcao_inv_g_inv}, está bem definida.

A função inversa possui as seguintes propriedades:
\begin{align}
  f^{-1}\left(f(x)\right) & = x, \qquad \forall~ x \in D\left(f\right),  \\[3mm]
  f\left(f^{-1}(x)\right) & = x, \qquad \forall~ x \in D\left(f^{-1}\right).
\end{align}

%------------------------------------------------------------------------------%
\begin{example}
  Determine se a função
  \begin{align*}
    f\colon \R & \to \R  \\
    x          & \to x^2
  \end{align*}
  é bijetiva.

  \solution
  A função $f$ não é bijetiva, pois,
  \[
    f(1) = f(-1) = 1.
  \]
\end{example}
%------------------------------------------------------------------------------%

%------------------------------------------------------------------------------%
\begin{example}
  Determine se a função
  \begin{align*}
    g\colon [0, \infty) & \to [0, \infty) \\
    x                   & \to  x^2
  \end{align*}
  é bijetiva.

  \solution
  A função $g$, com o domínio restringindo, é uma função bijetiva,
  cuja função inversa é dada por
  \begin{align*}
    f\colon [0, \infty) & \to [0, \infty) \\
    x                   & \to \sqrt{x}.
  \end{align*}
\end{example}
%------------------------------------------------------------------------------%

Como encontrar a função inversa de uma função bijetiva? Se $y=f(x)$ possuir
inversa, escreva $x=f(y)$ e isole $y$ na última equação.

%------------------------------------------------------------------------------%
\begin{example}
  Encontre, caso exista, a função inversa da Função
  \(
    f(x) = x^3 + 4
  \).
  Caso não exista, justifique. Suponha que essas funções tenham domínio e
  contradomínio todo o conjunto dos números reais.

  \solution
  Vamos verificar se $f$ é bijetora.

  Para verificar se $f$ é injetora, supomos que
  \[
    f(x_1) = f(x_2)
  \]
  para $x_1, x_2 \in \R$. Assim,
  \begin{align*}
    (x_1)^3+4     & = (x_2)^3+4     \\
    (x_1)^3       & = (x_2)^3       \\
    \sqrt[3]{x_1} & = \sqrt[3]{x_2} \\
    x_1           & = x_2
  \end{align*}
  Logo, $f$ é injetora.

  $f$ é sobrejetora, uma vez que seu conjunto imagem é todo o
  conjunto dos números reais, que é igual ao contradomínio.

  Para encontrarmos uma expressão para a função inversa, basta trocarmos $x$
  por $y$ na função e depois isolarmos $y$. Portanto,
  \begin{align*}
    x   & = y^3+4               \\
    y^3 & = x-4                 \\
    y   & = \sqrt[3]{x-4}
  \end{align*}
  Assim, $f$ possui inversa e
  \[
    f^{-1}(x)=\sqrt[3]{x-4}.
  \]
\end{example}
%------------------------------------------------------------------------------%

%------------------------------------------------------------------------------%
\begin{example}
  Encontre, caso exista, a função inversa da Função
  \(
    g(x)= x^4 - 7
  \).
  Caso não exista, justifique. Suponha que essas funções tenham domínio e
  contradomínio todo o conjunto dos números reais.

  \solution
  A função $g$ não é injetora, uma vez que
  \[
    g(1) = g(-1) = -6,
  \]
  logo não possui inversa.
\end{example}
%------------------------------------------------------------------------------%

%------------------------------------------------------------------------------%
\subsection{Gráfico das funções inversas}

O próximo objetivo é explorar as relações entre os gráficos de $f$ e de sua
função inversa $f^{-1}$. Se $(a, b)$ for um ponto no gráfico $y = f(x)$, então
$(b, a)$ será um ponto do gráfico de $y=f^{-1}(x)$. O efeito geométrico de
inverter as coordenadas de um ponto é refletir aquele ponto sobre a reta $y =
x$.

Em resumo, inverter as coordenadas de um ponto no gráfico de $f$ produz um
ponto no gráfico de $f^{-1}$. Analogamente inverter as coordenadas de um ponto
no gráfico de $f^{-1}$ produz um ponto no gráfico de  $f$ logo os gráficos de
$y = f(x)$ e $y = f^{-1}(x)$ são simétricos um do outro em relação a reta
$y=x$, conforme ilustra a Figura~\ref{fig:invespelho}.

%------------------------------------------------------------------------------%
\begin{figure}
  \centering
  \includegraphics[scale=0.6]{inversaespelho}
  \caption{Comparação do gráfico de $f(x)$ e $f^{-1}(x)$. Note que o gráfico de
  $f^{-1}(x)$ é o espelho do gráfico de $f(x)$, com relação a reta $y=x.$}
  \label{fig:invespelho}
\end{figure}
%------------------------------------------------------------------------------%

%------------------------------------------------------------------------------%